%% ---- 7. Discussão e Conclusão -----------------------------------------------
\section{Discussão e Conclusão}

Formulamos a explicação contrafactual de uma recomendação sequencial como a
busca de conjuntos de posições do histórico sob três objetivos (necessidade,
suficiência e tamanho), estimados por intervenções com fundos fixos, e a
resolvemos com algoritmos evolutivos multiobjetivo sobre máscaras binárias. A
proposta previa que uma busca multiobjetivo cobriria uma frente melhor que o
mesmo orçamento gasto em buscas escalarizadas, e apontava como risco que $N$ e
$S$ fossem quase redundantes. As duas coisas ocorreram, em regiões diferentes.
Nas explicações pequenas, $N$ e $S$ concordam só moderadamente, os algoritmos
baseados em dominância de Pareto são os melhores métodos, contra a busca
escalarizada, o guloso, o Shapley e o GECE, e manter $S$ como objetivo não
custa, em média, nada em necessidade e tamanho. Nas máscaras grandes, $N$ e
$S$ quase coincidem, e métodos que percorrem todos os tamanhos vencem o NSGA-II
original no MovieLens; só a variante semeada, posterior à comparação
planejada, vence ali. As conclusões dos métodos centrais se repetem em
instâncias novas e no Zvuk.

\paragraph{Limitações.}
(i)~Não há avaliação com usuários: mostramos que as frentes contêm explicações
pequenas, necessárias e às vezes suficientes, não que uma frente seja mais útil
a quem a lê do que uma resposta única. (ii)~Vários experimentos foram decididos
depois da comparação planejada ($\dagger$); a confirmação em dados novos cobre
os métodos centrais, mas não o NSGA-III, o Shapley e a variante só $(N,C)$, de
que depende o terceiro resultado. (iii)~A distribuição de fundo define o que
``suficiente'' significa, e trocá-la muda as explicações
(Apêndice~\ref{sec:ablacao}). (iv)~O ajuste foi desigual (12 configurações
para o NSGA-II, menos para os demais) e o MOEA/D só foi testado com pesos
uniformes; como os três algoritmos baseados em dominância diferem por menos de
$0{,}01$ de \HVp{}, mais ajuste poderia mudar a ordem entre eles. Outras
limitações estão no Apêndice~\ref{ap:limitacoes}.

\paragraph{Trabalho futuro.}
Uma semeadura que não custe nada nas explicações pequenas (por exemplo, só
com máscaras pequenas), operadores que usem a estrutura do problema (mutação
enviesada para os itens recentes), outras arquiteturas de recomendador e uma
avaliação com usuários.

\paragraph{Reprodutibilidade.}
O código, os arquivos de resultados usados nas tabelas e figuras, os roteiros
que as geram e os comandos de cada execução estão no repositório do projeto
(\url{https://github.com/jorgesilva2407/moo-seqrec-explanations}).
